% 04-methodology.tex -- draft. Facts below come from versions.env,
% .state/versoes-instaladas.txt, RUNBOOK.md and the pilot. Re-check versions
% against .state/versoes-instaladas.txt before submission.
\section{Testbed and Methodology}
\label{sec:method}

\subsection{Testbed}
\label{sec:testbed}

All experiments run on a single dedicated host with an Intel Core i7-12700K
(8 performance cores with two hardware threads each, plus 4 efficiency cores)
and 64\,GB of RAM. To reduce measurement noise, we set the frequency governor
to \texttt{performance}, disable turbo boost, and disable idle states deeper
than C1, following MeshInsight~\cite{zhu2023meshinsight}.

The cluster runs in four KVM virtual machines (Ubuntu 24.04, kernel
6.8.0-149) managed with libvirt. Every vCPU is pinned to a dedicated
physical thread of a performance core: the control-plane VM has 2 vCPUs and
6\,GB, and each of the three worker VMs has 4 vCPUs and 14\,GB. QEMU emulator
threads are pinned to a separate core pair, so no two VMs share a physical
core. Kubernetes is k3s v1.37.2, with Cilium 1.20.3 as the CNI in
every configuration, so that the baseline and the meshes share the same
underlying network. We install Istio 1.32.0, Linkerd
edge-26.11.1, and Cilium's service-mesh features from their upstream
releases with default proxy settings; all versions are pinned in the
artifact.

\paragraph{Mesh configuration}
Every mesh is configured with mutual authentication and encryption required.
For Istio (IS, IA4, IA7) we apply a mesh-wide \texttt{STRICT}
\texttt{PeerAuthentication}; for Linkerd, a default policy that only admits
authenticated mesh traffic; for Cilium, WireGuard encryption between nodes and
a \texttt{CiliumNetworkPolicy} with L7 HTTP rules on the application ports,
which activates the per-node Envoy. In IA7, each application namespace has its
own waypoint. Only the entry service of each application accepts plaintext
traffic from the load generator.

\paragraph{Placement}
Control-plane components (istiod, the Linkerd control plane, the Cilium
operator, CoreDNS, and Prometheus) run on the control-plane VM, so they do not
compete with the data plane. Application pods are assigned to the three workers
round-robin with a fixed order, so every configuration uses the same placement,
and each waypoint is pinned to a fixed worker. We remove the demonstration CPU
and memory limits and the liveness probes shipped with the benchmark manifests,
keeping requests and readiness probes: in a pilot, those limits throttled
individual services and caused restarts that would have dominated the
results. Mesh proxies keep their default resource settings; we verified that the
busiest sidecar was not CPU-throttled in our runs.

\paragraph{Applications}
We use three workloads. \emph{Echo} is a two-service chain of Fortio
servers (a relay forwarding each request to an echo server on another worker),
which isolates the per-hop cost of the data plane. \emph{Online Boutique}
(OB)~\cite{onlineboutique} has 11 services communicating over gRPC, with its
built-in load generator removed. \emph{Hotel Reservation}
(HR)~\cite{gan2019deathstarbench} has eight Go services that communicate over
gRPC, a Consul registry, and six MongoDB and three memcached instances. The k6 scripts emulate user sessions: for OB, browsing the home
page, viewing products, changing currency, and adding to the cart; for HR, 60\%
hotel search, 39\% recommendations, 0.5\% logins, and 0.5\% reservations.

\paragraph{Reproducible state}
Each configuration is deployed from the same base snapshot of all four VMs,
which contains the cluster and every container image already pulled, so no run
downloads images and every configuration starts from an identical state.
Automatic package updates are disabled inside the VMs.

\subsection{Load generation}
\label{sec:load}

The load generator, k6~\cite{k6}, runs on the host, pinned to the four
efficiency cores, which no VM uses. It issues requests at a constant arrival
rate (open loop), so a slow response does not delay subsequent requests and
tail latency is not hidden~\cite{tene2015latency}. At the highest rates we
used, the generator's cores were about 53\% utilized, and we track dropped
iterations to detect client saturation.

\paragraph{Saturation}
For each configuration and application, we find the maximum SLO-compliant
throughput by increasing the rate in steps of 25\%, each held for 60\,s, until
p99 latency exceeds the SLO (20\,ms for echo, 500\,ms for OB, 100\,ms for HR) or
more than 1\% of requests fail or are dropped, and then refine the boundary with four bisection steps
between the last passing and the first failing rate (a resolution of about
1.5\% of the rate).

\paragraph{Load levels}
All configurations must be compared at the same absolute request rate, and a
rate beyond a mesh's saturation would measure queue collapse rather than
overhead. We therefore set the load levels to 30\%, 60\%, and 90\% of the
\emph{lowest} saturation throughput among the six configurations for each
application. This yields 714, 1429, and 2143\,req/s for echo, 192, 385, and
578\,req/s for OB, and 259, 519, and 778\,req/s for HR.

\subsection{Experiment design and metrics}
\label{sec:design}

The main experiment crosses 6 configurations, 3 applications, and 3 load
levels, with 10 repetitions per cell, for 540 runs. Within each configuration,
the order of runs is randomized with a fixed seed. Each run has 60\,s of
warm-up, which is discarded, followed by 180\,s of measurement.

We collect: (i) latency percentiles (p50, p90, p99, p99.9), error rate, and
achieved throughput from k6; (ii) CPU and memory per container (application,
sidecar, ztunnel, waypoint, Cilium agent and Envoy, control plane) from
cAdvisor through Prometheus at a 5\,s interval; and (iii) package and core
energy from RAPL on the host. We subtract the idle power of the no-mesh
baseline from every configuration, so that each mesh's standing cost is
included, and report joules per request together with each mesh's idle power
above the baseline. Resource costs are also
normalized to millicores and MiB per 1{,}000\,req/s.

For profiling (RQ2), we repeat each configuration at the 60\% load level and
record call stacks with \texttt{perf} on the worker VMs, count context switches
and system calls per request with eBPF (bpftrace and BCC), and, for Cilium, read
the run time of each eBPF program through \texttt{kernel.bpf\_stats\_enabled}.

\paragraph{Statistics}
We report medians across repetitions with 95\% bootstrap confidence intervals.
We test for differences among configurations with the Kruskal--Wallis test and
compare each mesh with the baseline using the Mann--Whitney U test with Holm
correction, reporting Cliff's $\delta$ as the effect size. The coefficient of
variation of p99 across repetitions measures the predictability of each mesh.

\subsection{Validation against MeshInsight}
\label{sec:validation}

To check that a single-host testbed reproduces published behavior, we repeated
MeshInsight's comparison of the Istio sidecar against no mesh on Online
Boutique, with five repetitions per configuration. At 275\,req/s, the sidecar
increased cluster CPU usage by 129\%, within the published range of
42--163\%~\cite{zhu2023meshinsight}. Latency increased by 416\% (p50) and 440\%
(p99), above the published 27--269\%. At 551\,req/s the latency increase grew
further (p50 +1482\%), while CPU rose by 80\%. We attribute the higher latency
to two factors that we analyze in \cref{sec:rq1,sec:rq3}: a newer Istio and
Envoy, and the concentration of sidecar work on the node that hosts the
frontend, which reached 77\% utilization at 551\,req/s. The main experiment shows
the same effect: for the Istio sidecar on OB at 90\% load, the busiest worker
is 76\% utilized, against 43\% without a mesh.
